% part: intuitionistic-logic
% chapter: propositions-as-types

\documentclass[../../../include/open-logic-chapter]{subfiles}

\begin{document}

\olchapter{int}{pty}{بيانونه د ټايپونو په توګه}

\begin{editorial}
  دا د کري--هاورډ د مطابقت په باب د يوه فصل
  \emph{ډېره تجربوي} مسوده ده۔ نور وضاحت او د
  موضوع د اړتيا بيان غواړي، او ښايي تېروتنې او
  نيمګړتياوې ولري۔ د عادي کولو ثبوت بياکتنه او
  پراخول غواړي۔ د ضرب ټايپ دپاره نورې تفصيلي
  بېلګې پکار دي۔ د ځاے بدلولو او ساده کولو
  بدلونونه نۀ دي رانغښتل شوي۔ کله چې د ټايپ‌لرونکي
  لامبډا حساب مواد هم ورسره وي، ډېر به ئې مطلب
  څرګند شي، ځکه دلته اساساً د هغه پوهه له مخه
  فرض شوې ده۔ په ډېر احتياط ئې وکاروئ۔
  \textbf{د سرچينې د يادښت سپيناوی (OLCMP-211):}
  سرچينه د ضرب ټايپ بېلګې بالکل نشته بولي؛ خو د
  ثبوتونو په ترمونو د اړولو فصل د او د دوو برخو
  د بدلولو مثال لري چې ثبوت ترم ئې جوړه او دواړه
  پروجکشنونه کاروي۔ نو د مثال د بالکل نشتوالي پر
  ځاے د نورو تفصيلي بېلګو اړتيا دلته ويل شوې؛ د
  فصل د تجربوي حالت نور ټول اخطارونه ساتل شوي دي۔
\end{editorial}

\olimport{introduction}
%\olimport{sequent-natural-deduction}
\olimport{proof-terms}
\olimport{proofs-to-terms}
\olimport{terms-to-proofs}
\olimport{types}
\olimport{reduction}
\olimport{type-preservation}
\olimport{normalization}

\OLEndChapterHook

\end{document}
